\documentclass{article}
\begin{document}
\section{Alpha}\label{sec:a}
A typesetting system takes a description of a document and decides where every character goes. It breaks paragraphs into lines and stacks the lines into pages. See Section~\ref{sec:b} on page~\pageref{sec:b}.
\subsection{Beta}\label{sec:b}
Lines of a paragraph are chosen together rather than one at a time, so that the spacing is as even as the words allow and few hyphens appear. See \ref{sec:c} and \ref{sec:d}.
\subsubsection{Gamma}\label{sec:c}
A typesetting system takes a description of a document and decides where every character goes. It breaks paragraphs into lines and stacks the lines into pages.
\paragraph{Delta}\label{sec:d} Lines of a paragraph are chosen together rather than one at a time, so that the spacing is as even as the words allow and few hyphens appear. Back to \ref{sec:a}, page \pageref{sec:a}.
\subparagraph{Epsilon} A typesetting system takes a description of a document and decides where every character goes. It breaks paragraphs into lines and stacks the lines into pages.
\section{Zeta}\label{sec:e}
Lines of a paragraph are chosen together rather than one at a time, so that the spacing is as even as the words allow and few hyphens appear. This is Section~\ref{sec:e}, see page \pageref{sec:d}.
\subsection{Eta}
A typesetting system takes a description of a document and decides where every character goes. It breaks paragraphs into lines and stacks the lines into pages.
\end{document}
